\documentclass[12pt]{article}
\usepackage[utf8]{inputenc}
\usepackage[T1]{fontenc}
\usepackage{hyperref}
\title{Kaapor Sign Language Language Resources for K-12 Curriculum Development}
\author{Compiled by Claude Sonnet 5 (claude-sonnet-5) for Andr\'as Kornai}
\date{}
\begin{document}
\maketitle

\section{Introduction}
Kaapor Sign Language (ISO 639-3 uks; Glottocode urub1243) is a village sign language of the Ka'apor community in Maranh\~ao, Brazil. The names ``Urubu'' and ``Urubu-Kaapor'' are described as pejorative; the community name is Ka'apor. A 1968 count found 7 deaf people among about 500, and hearing community members commonly sign as well. Recent estimates cited in search results put current deaf users at roughly 7 to 13, with sources disagreeing on endangerment status. Documentation is sparse.

\section{Primary Educational Resources}
No curricula, primers or school programmes for this sign language were found in search results.

\section{Linguistic and Reference Materials}
\begin{itemize}
\item Wikipedia, Ka'apor Sign Language: \url{https://en.wikipedia.org/wiki/Ka%CA%BCapor_Sign_Language}
\item Wikipedia mirror, ISO 639:uks: \url{https://wikipedia.com/wiki/ISO_639:uks}
\item Endangered Languages Project entry: \url{https://www.endangeredlanguages.com/lang/7007?hl=de}
\item Culture in Crisis, ``Documentation and description of Urubu-Ka'apor Sign Language'' (ELDP-funded project with the University of S\~ao Paulo, ongoing): \url{https://cultureincrisis.org/projects/documentation-and-description-of-urub%C3%BA-kaapor-sign-language}
\item Rosetta Project entry: \url{https://rosettapanglossia.longnow.org/index.php?title=Uks}
\item CSULB South America language reference (PDF): \url{https://home.csulb.edu/~lemaster/South%20America/SA%20Brazil%20-%20Kaapor.pdf}
\end{itemize}

\section{Religious and Community Resources}
\begin{itemize}
\item Povos Ind\'igenas no Brasil, Ka'apor profile (about the people, in English): \url{https://pib.socioambiental.org/en/Povo:Ka'apor}
\item Povos Ind\'igenas no Brasil, Ka'apor profile (alternate): \url{https://indios.org.br/en/povo/kaapor}
\end{itemize}

\section{Digital and Audio Resources}
No video corpus or software was found in search results. The ELDP documentation project aims to record signers in everyday settings; consult its pages for any deposited archive.

\section{Curriculum Development Context}
Because the language is a shared village sign language, with deaf and hearing signers, any curriculum would depend on the community and on the documentation project's outputs. The spoken Ka'apor language (Wikipedia: \url{https://en.wikipedia.org/wiki/Ka%CA%BCapor_language}) is a separate language and relevant for bilingual planning.

\end{document}